% !TEX program = pdflatex
% Compilar: pdflatex mapa_supuestos && bibtex mapa_supuestos && pdflatex mapa_supuestos && pdflatex mapa_supuestos
\documentclass[11pt,a4paper,parskip=half]{scrartcl}
% ============================================================
% estilo.tex — preámbulo compartido por los tres documentos
% Se incluye con % ============================================================
% estilo.tex — preámbulo compartido por los tres documentos
% Se incluye con % ============================================================
% estilo.tex — preámbulo compartido por los tres documentos
% Se incluye con \input{estilo} después de \documentclass
% ============================================================

\usepackage[utf8]{inputenc}
\usepackage[T1]{fontenc}
\usepackage{lmodern}
\usepackage[spanish]{babel}

\usepackage[a4paper,margin=2.5cm]{geometry}
\usepackage{microtype}
\usepackage{enumitem}
\usepackage{booktabs}
\usepackage{longtable}
\usepackage{tabularx}
\usepackage{array}
\usepackage{colortbl}

\usepackage{xcolor}
\definecolor{ink}{RGB}{0,62,92}           % azul principal
\definecolor{inklight}{RGB}{92,132,153}   % azul apagado
\definecolor{rulegray}{RGB}{205,214,219}  % líneas divisorias
\definecolor{critred}{RGB}{158,38,32}     % prioridad crítica
\definecolor{amber}{RGB}{176,98,16}       % prioridad alta
\definecolor{slate}{RGB}{104,112,118}     % prioridad media
\definecolor{headrow}{RGB}{0,62,92}       % encabezado de tabla
\definecolor{softrow}{RGB}{242,246,248}   % fila alterna

\usepackage{natbib}   % debe cargarse antes que hyperref
\usepackage[colorlinks=true,linkcolor=ink,citecolor=inklight,urlcolor=inklight]{hyperref}

% ---------- títulos ----------
\addtokomafont{title}{\color{ink}\sffamily}
\addtokomafont{author}{\color{slate}\normalsize}
\addtokomafont{date}{\color{slate}\normalsize}
\addtokomafont{disposition}{\color{ink}\sffamily}
\setkomafont{section}{\color{ink}\sffamily\large\bfseries}
\setkomafont{subsection}{\color{ink}\sffamily\normalsize\bfseries}

% ---------- lista de campos alineada (reemplaza las cajas) ----------
\newlist{campos}{description}{1}
\setlist[campos]{
  labelwidth=2.9cm,
  leftmargin=3.2cm,
  labelsep=0.3cm,
  align=left,
  font=\normalfont\sffamily\bfseries\small\color{ink},
  itemsep=2.5pt,
  topsep=4pt,
  parsep=0pt,
  before=\small
}

% ---------- etiqueta de prioridad ----------
\newcommand{\prio}[2]{%
  {\setlength{\fboxsep}{3pt}%
   \colorbox{#1}{\textcolor{white}{\sffamily\bfseries\scriptsize\strut\,#2\,}}}%
}
\newcommand{\critica}{\prio{critred}{CRÍTICA}}
\newcommand{\alta}{\prio{amber}{ALTA}}
\newcommand{\media}{\prio{slate}{MEDIA}}

\newsavebox{\cajaetiqueta}
\newlength{\sepetiqueta}\setlength{\sepetiqueta}{10pt}

% ---------- encabezado de una entrada (pregunta o ficha) ----------
% #1 = clave corta (Q1 / P-04)   #2 = título   #3 = etiqueta derecha
\newcommand{\entrada}[3]{%
  \par\addvspace{14pt}%
  \noindent{\color{rulegray}\rule{\linewidth}{0.6pt}}%
  \par\addvspace{5pt}%
  \noindent
  \sbox\cajaetiqueta{#3}%
  \begin{minipage}[t]{\dimexpr\linewidth-\wd\cajaetiqueta-\sepetiqueta\relax}
    \raggedright\sffamily\bfseries\color{ink}#1\,\textcolor{rulegray}{\textbar}\, #2
  \end{minipage}\hspace{\sepetiqueta}\usebox\cajaetiqueta%
  \par\addvspace{3pt}%
}

% variante sin clave a la izquierda (para fichas bibliográficas)
\newcommand{\entradasimple}[2]{%
  \par\addvspace{14pt}%
  \noindent{\color{rulegray}\rule{\linewidth}{0.6pt}}%
  \par\addvspace{5pt}%
  \noindent
  \sbox\cajaetiqueta{#2}%
  \begin{minipage}[t]{\dimexpr\linewidth-\wd\cajaetiqueta-\sepetiqueta\relax}
    \raggedright\sffamily\bfseries\color{ink}#1
  \end{minipage}\hspace{\sepetiqueta}\usebox\cajaetiqueta%
  \par\addvspace{3pt}%
}

% permite cortar nombres largos de campo en \texttt sin salirse del margen
\newcommand{\campoTT}[2]{\texttt{#1\_}\allowbreak\texttt{#2}}

% ---------- nota destacada, sin caja pesada ----------
\newenvironment{nota}{%
  \par\addvspace{8pt}\small
  \begin{list}{}{%
    \setlength{\leftmargin}{10pt}\setlength{\rightmargin}{0pt}
    \setlength{\labelwidth}{0pt}\setlength{\itemindent}{0pt}
    \setlength{\listparindent}{0pt}}
  \item[]\color{black}%
}{%
  \end{list}\addvspace{6pt}%
}

% ---------- tablas ----------
\newcolumntype{L}[1]{>{\raggedright\arraybackslash}p{#1}}
\newcommand{\thead}[1]{\textcolor{white}{\sffamily\bfseries\small #1}}

\renewcommand{\arraystretch}{1.25}
\setlength{\tabcolsep}{5pt}
 después de \documentclass
% ============================================================

\usepackage[utf8]{inputenc}
\usepackage[T1]{fontenc}
\usepackage{lmodern}
\usepackage[spanish]{babel}

\usepackage[a4paper,margin=2.5cm]{geometry}
\usepackage{microtype}
\usepackage{enumitem}
\usepackage{booktabs}
\usepackage{longtable}
\usepackage{tabularx}
\usepackage{array}
\usepackage{colortbl}

\usepackage{xcolor}
\definecolor{ink}{RGB}{0,62,92}           % azul principal
\definecolor{inklight}{RGB}{92,132,153}   % azul apagado
\definecolor{rulegray}{RGB}{205,214,219}  % líneas divisorias
\definecolor{critred}{RGB}{158,38,32}     % prioridad crítica
\definecolor{amber}{RGB}{176,98,16}       % prioridad alta
\definecolor{slate}{RGB}{104,112,118}     % prioridad media
\definecolor{headrow}{RGB}{0,62,92}       % encabezado de tabla
\definecolor{softrow}{RGB}{242,246,248}   % fila alterna

\usepackage{natbib}   % debe cargarse antes que hyperref
\usepackage[colorlinks=true,linkcolor=ink,citecolor=inklight,urlcolor=inklight]{hyperref}

% ---------- títulos ----------
\addtokomafont{title}{\color{ink}\sffamily}
\addtokomafont{author}{\color{slate}\normalsize}
\addtokomafont{date}{\color{slate}\normalsize}
\addtokomafont{disposition}{\color{ink}\sffamily}
\setkomafont{section}{\color{ink}\sffamily\large\bfseries}
\setkomafont{subsection}{\color{ink}\sffamily\normalsize\bfseries}

% ---------- lista de campos alineada (reemplaza las cajas) ----------
\newlist{campos}{description}{1}
\setlist[campos]{
  labelwidth=2.9cm,
  leftmargin=3.2cm,
  labelsep=0.3cm,
  align=left,
  font=\normalfont\sffamily\bfseries\small\color{ink},
  itemsep=2.5pt,
  topsep=4pt,
  parsep=0pt,
  before=\small
}

% ---------- etiqueta de prioridad ----------
\newcommand{\prio}[2]{%
  {\setlength{\fboxsep}{3pt}%
   \colorbox{#1}{\textcolor{white}{\sffamily\bfseries\scriptsize\strut\,#2\,}}}%
}
\newcommand{\critica}{\prio{critred}{CRÍTICA}}
\newcommand{\alta}{\prio{amber}{ALTA}}
\newcommand{\media}{\prio{slate}{MEDIA}}

\newsavebox{\cajaetiqueta}
\newlength{\sepetiqueta}\setlength{\sepetiqueta}{10pt}

% ---------- encabezado de una entrada (pregunta o ficha) ----------
% #1 = clave corta (Q1 / P-04)   #2 = título   #3 = etiqueta derecha
\newcommand{\entrada}[3]{%
  \par\addvspace{14pt}%
  \noindent{\color{rulegray}\rule{\linewidth}{0.6pt}}%
  \par\addvspace{5pt}%
  \noindent
  \sbox\cajaetiqueta{#3}%
  \begin{minipage}[t]{\dimexpr\linewidth-\wd\cajaetiqueta-\sepetiqueta\relax}
    \raggedright\sffamily\bfseries\color{ink}#1\,\textcolor{rulegray}{\textbar}\, #2
  \end{minipage}\hspace{\sepetiqueta}\usebox\cajaetiqueta%
  \par\addvspace{3pt}%
}

% variante sin clave a la izquierda (para fichas bibliográficas)
\newcommand{\entradasimple}[2]{%
  \par\addvspace{14pt}%
  \noindent{\color{rulegray}\rule{\linewidth}{0.6pt}}%
  \par\addvspace{5pt}%
  \noindent
  \sbox\cajaetiqueta{#2}%
  \begin{minipage}[t]{\dimexpr\linewidth-\wd\cajaetiqueta-\sepetiqueta\relax}
    \raggedright\sffamily\bfseries\color{ink}#1
  \end{minipage}\hspace{\sepetiqueta}\usebox\cajaetiqueta%
  \par\addvspace{3pt}%
}

% permite cortar nombres largos de campo en \texttt sin salirse del margen
\newcommand{\campoTT}[2]{\texttt{#1\_}\allowbreak\texttt{#2}}

% ---------- nota destacada, sin caja pesada ----------
\newenvironment{nota}{%
  \par\addvspace{8pt}\small
  \begin{list}{}{%
    \setlength{\leftmargin}{10pt}\setlength{\rightmargin}{0pt}
    \setlength{\labelwidth}{0pt}\setlength{\itemindent}{0pt}
    \setlength{\listparindent}{0pt}}
  \item[]\color{black}%
}{%
  \end{list}\addvspace{6pt}%
}

% ---------- tablas ----------
\newcolumntype{L}[1]{>{\raggedright\arraybackslash}p{#1}}
\newcommand{\thead}[1]{\textcolor{white}{\sffamily\bfseries\small #1}}

\renewcommand{\arraystretch}{1.25}
\setlength{\tabcolsep}{5pt}
 después de \documentclass
% ============================================================

\usepackage[utf8]{inputenc}
\usepackage[T1]{fontenc}
\usepackage{lmodern}
\usepackage[spanish]{babel}

\usepackage[a4paper,margin=2.5cm]{geometry}
\usepackage{microtype}
\usepackage{enumitem}
\usepackage{booktabs}
\usepackage{longtable}
\usepackage{tabularx}
\usepackage{array}
\usepackage{colortbl}

\usepackage{xcolor}
\definecolor{ink}{RGB}{0,62,92}           % azul principal
\definecolor{inklight}{RGB}{92,132,153}   % azul apagado
\definecolor{rulegray}{RGB}{205,214,219}  % líneas divisorias
\definecolor{critred}{RGB}{158,38,32}     % prioridad crítica
\definecolor{amber}{RGB}{176,98,16}       % prioridad alta
\definecolor{slate}{RGB}{104,112,118}     % prioridad media
\definecolor{headrow}{RGB}{0,62,92}       % encabezado de tabla
\definecolor{softrow}{RGB}{242,246,248}   % fila alterna

\usepackage{natbib}   % debe cargarse antes que hyperref
\usepackage[colorlinks=true,linkcolor=ink,citecolor=inklight,urlcolor=inklight]{hyperref}

% ---------- títulos ----------
\addtokomafont{title}{\color{ink}\sffamily}
\addtokomafont{author}{\color{slate}\normalsize}
\addtokomafont{date}{\color{slate}\normalsize}
\addtokomafont{disposition}{\color{ink}\sffamily}
\setkomafont{section}{\color{ink}\sffamily\large\bfseries}
\setkomafont{subsection}{\color{ink}\sffamily\normalsize\bfseries}

% ---------- lista de campos alineada (reemplaza las cajas) ----------
\newlist{campos}{description}{1}
\setlist[campos]{
  labelwidth=2.9cm,
  leftmargin=3.2cm,
  labelsep=0.3cm,
  align=left,
  font=\normalfont\sffamily\bfseries\small\color{ink},
  itemsep=2.5pt,
  topsep=4pt,
  parsep=0pt,
  before=\small
}

% ---------- etiqueta de prioridad ----------
\newcommand{\prio}[2]{%
  {\setlength{\fboxsep}{3pt}%
   \colorbox{#1}{\textcolor{white}{\sffamily\bfseries\scriptsize\strut\,#2\,}}}%
}
\newcommand{\critica}{\prio{critred}{CRÍTICA}}
\newcommand{\alta}{\prio{amber}{ALTA}}
\newcommand{\media}{\prio{slate}{MEDIA}}

\newsavebox{\cajaetiqueta}
\newlength{\sepetiqueta}\setlength{\sepetiqueta}{10pt}

% ---------- encabezado de una entrada (pregunta o ficha) ----------
% #1 = clave corta (Q1 / P-04)   #2 = título   #3 = etiqueta derecha
\newcommand{\entrada}[3]{%
  \par\addvspace{14pt}%
  \noindent{\color{rulegray}\rule{\linewidth}{0.6pt}}%
  \par\addvspace{5pt}%
  \noindent
  \sbox\cajaetiqueta{#3}%
  \begin{minipage}[t]{\dimexpr\linewidth-\wd\cajaetiqueta-\sepetiqueta\relax}
    \raggedright\sffamily\bfseries\color{ink}#1\,\textcolor{rulegray}{\textbar}\, #2
  \end{minipage}\hspace{\sepetiqueta}\usebox\cajaetiqueta%
  \par\addvspace{3pt}%
}

% variante sin clave a la izquierda (para fichas bibliográficas)
\newcommand{\entradasimple}[2]{%
  \par\addvspace{14pt}%
  \noindent{\color{rulegray}\rule{\linewidth}{0.6pt}}%
  \par\addvspace{5pt}%
  \noindent
  \sbox\cajaetiqueta{#2}%
  \begin{minipage}[t]{\dimexpr\linewidth-\wd\cajaetiqueta-\sepetiqueta\relax}
    \raggedright\sffamily\bfseries\color{ink}#1
  \end{minipage}\hspace{\sepetiqueta}\usebox\cajaetiqueta%
  \par\addvspace{3pt}%
}

% permite cortar nombres largos de campo en \texttt sin salirse del margen
\newcommand{\campoTT}[2]{\texttt{#1\_}\allowbreak\texttt{#2}}

% ---------- nota destacada, sin caja pesada ----------
\newenvironment{nota}{%
  \par\addvspace{8pt}\small
  \begin{list}{}{%
    \setlength{\leftmargin}{10pt}\setlength{\rightmargin}{0pt}
    \setlength{\labelwidth}{0pt}\setlength{\itemindent}{0pt}
    \setlength{\listparindent}{0pt}}
  \item[]\color{black}%
}{%
  \end{list}\addvspace{6pt}%
}

% ---------- tablas ----------
\newcolumntype{L}[1]{>{\raggedright\arraybackslash}p{#1}}
\newcommand{\thead}[1]{\textcolor{white}{\sffamily\bfseries\small #1}}

\renewcommand{\arraystretch}{1.25}
\setlength{\tabcolsep}{5pt}

\usepackage{pdflscape}
\bibliographystyle{apalike}

\title{Mapa de supuestos}
\subject{\normalsize ICS2122 — Taller de Investigación Operativa (Capstone)}
\subtitle{\normalsize Enunciado $\times$ literatura $\times$ decisión propia\\[2pt]
Planificación de la operación de buses eléctricos, Red Metropolitana de Movilidad}
\author{}
\date{}

\begin{document}
\maketitle
\vspace{-2.2\baselineskip}

\begin{nota}
\textbf{\textcolor{ink}{Cómo leer la tabla.}} Cada fila cruza tres cosas: qué dice —o qué calla— el enunciado sobre una dimensión del problema; cómo la trata la literatura verificada en \texttt{bibliografia.tex} (referencias \textbf{P-XX}); y qué decidimos nosotros, con la pregunta \textbf{QXX} de \texttt{preguntas\_profesor.tex} que respalda esa decisión. Cuando la celda dice \textit{«resuelto»}, el enunciado o los datos ya fijan el supuesto y no hace falta preguntar.
\end{nota}

\begin{landscape}

\section*{1\quad Supuestos que el enunciado sí plantea}
\addcontentsline{toc}{section}{Supuestos que el enunciado sí plantea}

\begin{longtable}{@{}L{4.4cm}L{8.4cm}L{8.4cm}@{}}
\rowcolor{headrow}
\thead{Dimensión} & \thead{Cómo lo modela la literatura} & \thead{Qué decidimos nosotros} \\
\endfirsthead
\rowcolor{headrow}
\thead{Dimensión} & \thead{Cómo lo modela la literatura} & \thead{Qué decidimos nosotros} \\
\endhead
\midrule

\textbf{Cobertura exacta de viajes.} Cada viaje se cubre una vez, sin duplicidad ni omisión. &
VSP/MDVSP clásico como problema de \textit{set partitioning} o \textit{covering}: P-01, P-02. &
Restricción de cobertura exacta sobre el GTFS de un día tipo. \textcolor{slate}{Q3 fija cuál día.} \\
\addlinespace
\rowcolor{softrow}
\textbf{Encadenamiento factible.} Un bus solo sigue a otro viaje si alcanza a terminar, desplazarse y llegar a horario. &
Arcos factibles en red \textit{connection-based} o \textit{time-space}: P-01, P-02. &
\textit{Time-space network} agregado \citep{kliewer_timespacemdvsp_2006}. \textcolor{slate}{Resuelto: el enunciado deja la metodología libre y exige justificarla.} \\
\addlinespace
\textbf{Autonomía de batería.} El bus consume energía al operar y no puede quedar sin carga. &
SoC como dimensión de estado: P-08, P-09, P-10. P-21 muestra que el consumo real depende de ocupación y conductor, no solo de distancia. &
Consumo lineal a 1{,}4~kWh/km, con SoC entre \texttt{min\_soc} y \texttt{max\_soc}. Se documenta la brecha frente a 1{,}006~kWh/km medido en Santiago. \textbf{Q5.} \\
\addlinespace
\rowcolor{softrow}
\textbf{Cargadores finitos.} Número limitado de cargadores por electroterminal. &
Cupos discretos de carga simultánea: P-11, P-12, P-13. &
\textbf{No hay dato:} \texttt{num\_chargers} está documentado pero ausente, y \campoTT{max\_simultaneous}{charging} iguala la capacidad de estacionamiento, así que no restringe. \textbf{Q2.} \\
\addlinespace
\textbf{Capacidad eléctrica restringida.} Distinta del número de cargadores. &
Potencia agregada en MW: P-23 la dimensiona a escala de sistema; P-06 aporta la perspectiva de red eléctrica. &
\textbf{No hay dato:} \texttt{transformer\_mw} documentado pero ausente. Proponemos 6~MW por electroterminal calibrado con P-20, con sensibilidad. \textbf{Q2.} \\
\addlinespace
\rowcolor{softrow}
\textbf{Deadhead con costo.} Los desplazamientos sin pasajeros generan tiempo y costo. &
Ubicación de electroterminales como dato frente a decisión: P-14 la optimiza; P-15 advierte que electrificar puede obligar a rediseñar toda la cadena de planificación. &
Ubicaciones fijas de \texttt{depots.csv}. Las distancias de \textit{deadhead} no vienen en ningún archivo y hay que calcularlas. \textbf{Q1.} \\
\addlinespace
\textbf{Trade-off energía vs.\ deadhead.} Una solución barata en energía puede ser cara en kilómetros vacíos. &
P-04 reporta que la literatura suele tratar flota, \textit{deadhead} y energía como objetivos separados. &
Función objetivo única aditiva, con análisis posterior del intercambio. \textbf{Q7.} \\
\addlinespace
\rowcolor{softrow}
\textbf{Flota finita.} La operación usa un número finito de vehículos. &
Flota fija: P-09. Flota como decisión de inversión (no aplica a un horizonte de un día): P-17. &
\texttt{fleet\_size = 1200} como cota superior. \textcolor{slate}{Resuelto: el enunciado dice que la programación influye en «la cantidad de buses necesarios».} \\
\addlinespace
\textbf{Metodología libre.} Exacto, heurístico, metaheurístico, híbrido o descomposición. &
Exactos: P-02, P-08, P-12, P-16. Heurísticos: P-07, P-10. Descomposición: P-17. &
Dos etapas: exacto en instancia reducida, luego ALNS extendida para la instancia completa. \textbf{Q12}, \textbf{Q13.} \\
\bottomrule
\end{longtable}

\clearpage
\section*{2\quad Supuestos que el enunciado no plantea y hay que fijar}
\addcontentsline{toc}{section}{Supuestos que el enunciado no plantea}

\begin{longtable}{@{}L{4.4cm}L{8.4cm}L{8.4cm}@{}}
\rowcolor{headrow}
\thead{Dimensión} & \thead{Cómo lo modela la literatura} & \thead{Qué decidimos nosotros} \\
\endfirsthead
\rowcolor{headrow}
\thead{Dimensión} & \thead{Cómo lo modela la literatura} & \thead{Qué decidimos nosotros} \\
\endhead
\midrule

\textbf{Tabla de viajes.} Distancia, duración, energía y \textit{deadhead} por viaje. &
No aplica: es un hallazgo de auditoría de datos, no un supuesto de la literatura. &
17 de los 44 campos documentados no existen en ningún archivo. Se construyen desde el GTFS crudo. \textbf{Q1 (bloqueante).} \\
\addlinespace
\rowcolor{softrow}
\textbf{Día tipo y escala.} El enunciado habla de «operación diaria» sin decir cuál día. &
No aplica. &
Día laboral (\texttt{service\_id = L}) sobre la red completa; todo recorte se declara. \textbf{Q3.} \\
\addlinespace
\textbf{Curva de carga.} ¿Lineal o no lineal? &
P-08 compara explícitamente ambos modelos, con depreciación de batería por profundidad de descarga. &
Carga a potencia constante, consistente con modelo lineal. \textbf{Q5.} \\
\addlinespace
\rowcolor{softrow}
\textbf{Fin del día.} ¿Retorno al electroterminal de origen? ¿Carga mínima al cierre? &
En el VSP clásico el retorno es a \emph{algún} depósito, no al de origen: P-01. &
Termina en cualquier electroterminal con cupo, respetando el SoC mínimo al cierre. \textbf{Q6.} \\
\addlinespace
\textbf{Incertidumbre.} Los datos son valores puntuales, sin distribuciones. &
P-04 identifica la robustez como el vacío más citado del área; P-18 la aborda con restricciones probabilísticas; P-05 confirma que sigue abierta en 2025. &
Determinista con sensibilidad ex-post ($\pm$15\,\% consumo, $\pm$10\,\% \textit{deadhead}). \textbf{Q10.} \\
\addlinespace
\rowcolor{softrow}
\textbf{Baseline y criterio de éxito.} No se indica contra qué comparar. &
P-21 y P-22 comparan contra operación diésel, comparación imposible aquí por falta de datos. &
Heurística \textit{greedy} propia más \textit{gap} contra la relajación lineal. \textbf{Q9.} \\
\addlinespace
\textbf{Herramientas de software.} No se mencionan solvers. &
P-16 alcanza 101 viajes a optimalidad con \textit{branch-and-price} sobre solver comercial. &
Se asume Gurobi académico; si no está, se reducen las instancias exactas. \textbf{Q11.} \\
\addlinespace
\rowcolor{softrow}
\textbf{Escala real vs.\ datos.} El enunciado describe la RM completa; los datos dan 1.200 buses y 5 electroterminales. &
RED tenía 6.756 buses y 380 rutas en 2018 según P-20. &
\texttt{01-datos} como universo cerrado; las fuentes reales solo para discutir realismo. \textbf{Q4.} \\
\addlinespace
\textbf{Exclusiones.} Intercambio de baterías y turnos de conductores. &
P-10 modela intercambio de baterías; P-03 ubica el \textit{crew scheduling} como problema hermano. &
Ambos fuera de alcance, declarado explícitamente. \textbf{Q8.} \\
\bottomrule
\end{longtable}

\end{landscape}

\clearpage
\renewcommand{\refname}{Referencias citadas}
\bibliography{referencias}

\end{document}
